% GENERATED FILE. Edit canonical lessons, then run npm run generate:books.
% canonical-source-hash: fnv1a64:080c26ec43dce516
% canonical-lessons: FA-C309-sokhangu, FA-C309-bayaniye, FA-C309-ettelaiye, FA-C309-elam-kardan, FA-C309-ezhar-kardan

\chapter{Who said what}
\label{ch:fa-who-said-what}
\hlchaptermodality{309}

\begin{chapteropening}
\textbf{By the end of this chapter:} \emph{I can say who made a statement: a spokesperson, a statement, a notice, announcing and stating.}

Complete the last lesson of chapter 309: I can say who made a statement: a spokesperson, a statement, a notice, announcing and stating.
\end{chapteropening}

\section[\texorpdfstring{\fa{سخنگو} (sokhangu)}{سخنگو (sokhangu)}]{\fa{سخنگو} (sokhangu) — a spokesperson}

\label{lesson:FA-C309-sokhangu}



\begin{quote}
Before the new one: say the Persian for a headline, then the Persian for the news.
\end{quote}

\subsection*{You'll want to know: \fa{سخنگو}}
\textbf{\fa{سخنگو}} (\emph{sokhangu}) — \textquotedblleft{}a spokesperson\textquotedblright{}.

\begin{quote}
\fa{سخنگوی} \fa{پلیس} \fa{با} \fa{خبرنگاران} \fa{صحبت} \fa{کرد}.
\end{quote}

\begin{quote}
\emph{Sokhangu-ye polis bâ khabarnegârân sohbat kard.}
\end{quote}

\begin{quote}
The police spokesperson spoke to the journalists.
\end{quote}

\emph{sokhan} \textquotedblleft{}speech\textquotedblright{} + \emph{gu} \textquotedblleft{}saying\textquotedblright{}, the present stem of \emph{goftan} \textquotedblleft{}to say\textquotedblright{}.

\begin{cousinweb}
\emph{sokhangu} joins \emph{sokhan} \textquotedblleft{}speech\textquotedblright{} and \emph{gu-}, the present stem of \emph{goftan} \textquotedblleft{}to say\textquotedblright{}: the speech-sayer. Both parts are native Persian; EtymDB lists no entry.
\end{cousinweb}

\begin{grammarlens}[title={What you've built}]
One more word for telling, explaining and describing at B1.
\end{grammarlens}

\subsection*{Your turn}
\begin{itemize}
\raggedright
  \item \emph{Say it:} \emph{sokhangu}
  \item \emph{Say it:} \emph{sokhangu}, once more
  \item \emph{Say it:} a sentence of your own with \emph{sokhangu}
  \item \emph{Recall:} say the Persian for a headline, then the Persian for the news, then say \emph{sokhangu} again
\end{itemize}

\begin{tcolorbox}[breakable,colback=teal!4,colframe=teal!35!black,title={Before you move on}]
What does \fa{سخنگو} mean? (\textquotedblleft{}A spokesperson\textquotedblright{}.) Say it once more.
\end{tcolorbox}
\section[\texorpdfstring{\fa{بیانیه} (bayâniye)}{بیانیه (bayâniye)}]{\fa{بیانیه} (bayâniye) — a statement, a communiqué}

\label{lesson:FA-C309-bayaniye}



\begin{quote}
Before the new one: say the Persian for the news, then the Persian for a spokesperson.
\end{quote}

\subsection*{You'll want to know: \fa{بیانیه}}
\textbf{\fa{بیانیه}} (\emph{bayâniye}) — \textquotedblleft{}a statement, a communiqué\textquotedblright{}.

\begin{quote}
\fa{شرکت} \fa{دیروز} \fa{یک} \fa{بیانیه} \fa{کوتاه} \fa{داد}.
\end{quote}

\begin{quote}
\emph{Sherkat diruz yek bayâniye-ye kutâh dâd.}
\end{quote}

\begin{quote}
Yesterday the company issued a short statement.
\end{quote}

Formal; \emph{bayâniye dâdan} \textquotedblleft{}to issue a statement\textquotedblright{}. Built on \emph{bayân} \textquotedblleft{}expression\textquotedblright{}.

\begin{cousinweb}
\emph{bayâniye} is built on \emph{bayân} \textquotedblleft{}expression, statement\textquotedblright{}, which is borrowed from Arabic \emph{bayān} (EtymDB lists both words), with the noun ending \emph{-iye}.
\end{cousinweb}

\begin{grammarlens}[title={What you've built}]
One more word for telling, explaining and describing at B1.
\end{grammarlens}

\subsection*{Your turn}
\begin{itemize}
\raggedright
  \item \emph{Say it:} \emph{bayâniye}
  \item \emph{Say it:} \emph{bayâniye}, once more
  \item \emph{Say it:} a sentence of your own with \emph{bayâniye}
  \item \emph{Recall:} say the Persian for the news, then the Persian for a spokesperson, then say \emph{bayâniye} again
\end{itemize}

\begin{tcolorbox}[breakable,colback=teal!4,colframe=teal!35!black,title={Before you move on}]
What does \fa{بیانیه} mean? (\textquotedblleft{}A statement, a communiqué\textquotedblright{}.) Say it once more.
\end{tcolorbox}
\section[\texorpdfstring{\fa{اطلاعیه} (ettelâ'iye)}{اطلاعیه (ettelâ'iye)}]{\fa{اطلاعیه} (ettelâ'iye) — a notice, an announcement}

\label{lesson:FA-C309-ettelaiye}



\begin{quote}
Before the new one: say the Persian for a spokesperson, then the Persian for a statement.
\end{quote}

\subsection*{You'll want to know: \fa{اطلاعیه}}
\textbf{\fa{اطلاعیه}} (\emph{ettelâ'iye}) — \textquotedblleft{}a notice, an announcement\textquotedblright{}.

\begin{quote}
\fa{روی} \fa{در} \fa{مدرسه} \fa{یک} \fa{اطلاعیه} \fa{بود}.
\end{quote}

\begin{quote}
\emph{Ru-ye dar-e madrese yek ettelâ'iye bud.}
\end{quote}

\begin{quote}
There was a notice on the school door.
\end{quote}

Built on \emph{ettelâ'} \textquotedblleft{}information\textquotedblright{}. Also a printed notice in an office or a mosque.

\begin{cousinweb}
\emph{ettelâ'iye} is built on \emph{ettelâ'} \textquotedblleft{}information, knowledge\textquotedblright{}, which is usually traced to Arabic \emph{iṭṭilā'}; EtymDB lists the Persian word with no parent.
\end{cousinweb}

\begin{grammarlens}[title={What you've built}]
One more word for telling, explaining and describing at B1.
\end{grammarlens}

\subsection*{Your turn}
\begin{itemize}
\raggedright
  \item \emph{Say it:} \emph{ettelâ'iye}
  \item \emph{Say it:} \emph{ettelâ'iye}, once more
  \item \emph{Say it:} a sentence of your own with \emph{ettelâ'iye}
  \item \emph{Recall:} say the Persian for a spokesperson, then the Persian for a statement, then say \emph{ettelâ'iye} again
\end{itemize}

\begin{tcolorbox}[breakable,colback=teal!4,colframe=teal!35!black,title={Before you move on}]
What does \fa{اطلاعیه} mean? (\textquotedblleft{}A notice, an announcement\textquotedblright{}.) Say it once more.
\end{tcolorbox}
\section[\texorpdfstring{\fa{اعلام} \fa{کردن} (e'lâm kardan)}{اعلام کردن (e'lâm kardan)}]{\fa{اعلام} \fa{کردن} (e'lâm kardan) — to announce}

\label{lesson:FA-C309-elam-kardan}



\begin{quote}
Before the new one: say the Persian for a statement, then the Persian for a notice.
\end{quote}

\subsection*{You'll want to know: \fa{اعلام} \fa{کردن}}
\textbf{\fa{اعلام} \fa{کردن}} (\emph{e'lâm kardan}) — \textquotedblleft{}to announce\textquotedblright{}.

\begin{quote}
\fa{مدرسه} \fa{اعلام} \fa{کرد} \fa{که} \fa{فردا} \fa{تعطیل} \fa{است}.
\end{quote}

\begin{quote}
\emph{Madrese e'lâm kard ke fardâ ta'til ast.}
\end{quote}

\begin{quote}
The school announced that it would be closed tomorrow.
\end{quote}

Present stem \emph{kon-}: \emph{e'lâm mikonad}. The passive \emph{e'lâm shod} is \textquotedblleft{}it was announced\textquotedblright{}.

\begin{cousinweb}
\emph{e'lâm} is borrowed from Arabic \emph{i'lām}, \textquotedblleft{}notice, instruction\textquotedblright{} (EtymDB). \emph{kardan} \textquotedblleft{}to do\textquotedblright{} is inherited Persian.
\end{cousinweb}

\begin{grammarlens}[title={What you've built}]
One more word for telling, explaining and describing at B1.
\end{grammarlens}

\subsection*{Your turn}
\begin{itemize}
\raggedright
  \item \emph{Say it:} \emph{e'lâm kardan}
  \item \emph{Say it:} \emph{e'lâm kardan}, once more
  \item \emph{Say it:} a sentence of your own with \emph{e'lâm kardan}
  \item \emph{Recall:} say the Persian for a statement, then the Persian for a notice, then say \emph{e'lâm kardan} again
\end{itemize}

\begin{tcolorbox}[breakable,colback=teal!4,colframe=teal!35!black,title={Before you move on}]
What does \fa{اعلام} \fa{کردن} mean? (\textquotedblleft{}To announce\textquotedblright{}.) Say it once more.
\end{tcolorbox}
\section[\texorpdfstring{\fa{اظهار} \fa{کردن} (ezhâr kardan)}{اظهار کردن (ezhâr kardan)}]{\fa{اظهار} \fa{کردن} (ezhâr kardan) — to state, to declare}

\label{lesson:FA-C309-ezhar-kardan}



\begin{quote}
Before the new one: say the Persian for a notice, then the Persian for to announce.
\end{quote}

\subsection*{You'll want to know: \fa{اظهار} \fa{کردن}}
\textbf{\fa{اظهار} \fa{کردن}} (\emph{ezhâr kardan}) — \textquotedblleft{}to state, to declare\textquotedblright{}.

\begin{quote}
\fa{مرد} \fa{اظهار} \fa{کرد} \fa{که} \fa{چیزی} \fa{ندیده} \fa{است}.
\end{quote}

\begin{quote}
\emph{Mard ezhâr kard ke chizi nadide ast.}
\end{quote}

\begin{quote}
The man stated that he had not seen anything.
\end{quote}

Formal, mostly in reports. Present \emph{ezhâr mikonad}. \emph{ezhâr-e nazar} \textquotedblleft{}a comment\textquotedblright{}.

\begin{cousinweb}
\emph{ezhâr} is borrowed from Arabic \emph{iẓhār} (EtymDB), usually glossed \textquotedblleft{}a showing, a making plain\textquotedblright{}. \emph{kardan} \textquotedblleft{}to do\textquotedblright{} is inherited Persian.
\end{cousinweb}

\begin{grammarlens}[title={What you've built}]
One more word for telling, explaining and describing at B1.
\end{grammarlens}

\subsection*{Your turn}
\begin{itemize}
\raggedright
  \item \emph{Say it:} \emph{ezhâr kardan}
  \item \emph{Say it:} \emph{ezhâr kardan}, once more
  \item \emph{Say it:} a sentence of your own with \emph{ezhâr kardan}
  \item \emph{Recall:} say the Persian for a notice, then the Persian for to announce, then say \emph{ezhâr kardan} again
\end{itemize}

\begin{tcolorbox}[breakable,colback=teal!4,colframe=teal!35!black,title={Before you move on}]
What does \fa{اظهار} \fa{کردن} mean? (\textquotedblleft{}To state, to declare\textquotedblright{}.) Say it once more.
\end{tcolorbox}
